% --- Archon physics formalization source begin ---
% archon:physics
% archon:covers PhyXMiniProblems/problem_phyx_mini_0290.lean
% archon:source-report reports/phyx_mini/problem_phyx_mini_0290.source.json

\chapter{Physics problem phyx\_mini\_0290}
\label{ch:PhyXMiniProblems_problem_phyx_mini_0290}

\paragraph{Problem source.}
\#\# Physical scenario

A sinusoidal transverse wave of wavelength \$20\$ cm travels along a string in the positive direction of an \$x\$ axis. The displacement \$y\$ of the string particle at \$x = 0\$ is given in figure as a function of time \$t\$. The scale of the vertical axis is set by \$y\_s = 4.0\$ cm. The wave equation is to be in the form \$y(x, t) = y\_m \textbackslash{}sin(kx \textbackslash{}pm \textbackslash{}omega t + \textbackslash{}phi)\$.

\#\# Question

What is the transverse velocity of the particle at \$x = 0\$ when \$t = 5.0\$ s?

\#\# Answer choices

- A: A: -1.5 cm/s

- B: B: -2.0 cm/s

- C: C: -3.0 cm/s

- D: D: -2.5 cm/s

\#\# Figure caption (auxiliary; use the image as primary evidence)

The image is a graph showing a sine wave. Here are the details:

- The horizontal axis is labeled as \textbackslash{}( t \textbackslash{}, (s) \textbackslash{}), representing time in seconds.
- The vertical axis is labeled as \textbackslash{}( y \textbackslash{}, (cm) \textbackslash{}), representing a displacement or position in centimeters.
- The wave oscillates between \textbackslash{}( y\_s \textbackslash{}) and \textbackslash{}( -y\_s \textbackslash{}) on the vertical axis.
- The horizontal axis shows values ranging from 0 to 10.
- The graph is a sinusoidal wave that begins at 0, reaches a peak at \textbackslash{}( y\_s \textbackslash{}), descends to \textbackslash{}( -y\_s \textbackslash{}), and then ascends back to a peak again as time progresses.
- The grid lines help show the periodic nature of the wave.
- The thickness of the wave line is relatively bold, highlighting the waveform against the grid.

\#\# Dataset metadata

Wave Properties; Physical Model Grounding Reasoning, Numerical Reasoning

\paragraph{Recorded answer/context.}
D: D: -2.5 cm/s

\paragraph{Figure/image path.}
/root/proposal\_for\_physic/science-mango/phyx\_mini\_run/phyx\_data/test\_image/290.png

\paragraph{Formalization target.}
create a compiling Lean file with sorry bodies at `PhyXMiniProblems/problem\_phyx\_mini\_0290.lean`. The Lean declarations must preserve the physical quantities, dimensions or dimensional roles, figure labels, governing-law hypotheses, and final relation expressed by this problem.
Use Mathlib/Physlib names found through LeanExplore where available. If a physics API is missing, introduce faithful local abstractions rather than scalar placeholder aliases.

\begin{theorem}[Physics formalization target]
\label{thm:physics:phyx_mini_0290:target}
\lean{PhyXMiniProblems.ProblemPhyXMini0290.problem_phyx_mini_0290}
\uses{def:physics:phyx-mini-0290:phyxminiproblems-problemphyxmini0290-transversevelocityreadout, def:physics:phyx-mini-0290:phyxminiproblems-problemphyxmini0290-travelingstringwavesetup, def:physics:phyx-mini-0290:phyxminiproblems-problemphyxmini0290-matchesstringwaveproblemdata, def:physics:phyx-mini-0290:phyxminiproblems-problemphyxmini0290-matchesprimarywavegraph, def:physics:phyx-mini-0290:phyxminiproblems-problemphyxmini0290-hasphysicaltravelingwaveparameters, def:physics:phyx-mini-0290:phyxminiproblems-problemphyxmini0290-satisfiessinusoidaltravelingwavelaws, def:physics:phyx-mini-0290:phyxminiproblems-problemphyxmini0290-answerchoice, def:physics:phyx-mini-0290:phyxminiproblems-problemphyxmini0290-displayedvelocityincentimeterspersecond, def:physics:phyx-mini-0290:phyxminiproblems-problemphyxmini0290-recordeddatasetanswer, def:physics:phyx-mini-0290:phyxminiproblems-problemphyxmini0290-matchesdisplayedvelocity}
The assigned autoformalize agent should translate this physics problem into Lean declarations in the covered file, with theorem and lemma proof bodies written as `by sorry`.
\end{theorem}
\begin{proof}
This is an autoformalization task, not a proof task. Produce statements that can later be proved by the physics prover without weakening the physical model.
\end{proof}

% --- Archon declaration topology begin ---
\section{Declaration topology}
\begin{definition}[Amplitude Quantity]
  \label{def:physics:phyx-mini-0290:phyxminiproblems-problemphyxmini0290-amplitudequantity}
  \lean{PhyXMiniProblems.ProblemPhyXMini0290.AmplitudeQuantity}
  A nonnegative transverse amplitude
\end{definition}
\begin{proof}
  This is the declared model definition of Amplitude Quantity.
\end{proof}
\begin{definition}[Wavelength Quantity]
  \label{def:physics:phyx-mini-0290:phyxminiproblems-problemphyxmini0290-wavelengthquantity}
  \lean{PhyXMiniProblems.ProblemPhyXMini0290.WavelengthQuantity}
  A nonnegative physical wavelength
\end{definition}
\begin{proof}
  This is the declared model definition of Wavelength Quantity.
\end{proof}
\begin{definition}[Signed Length Quantity]
  \label{def:physics:phyx-mini-0290:phyxminiproblems-problemphyxmini0290-signedlengthquantity}
  \lean{PhyXMiniProblems.ProblemPhyXMini0290.SignedLengthQuantity}
  A signed axial coordinate or transverse displacement
\end{definition}
\begin{proof}
  This is the declared model definition of Signed Length Quantity.
\end{proof}
\begin{definition}[Time Quantity]
  \label{def:physics:phyx-mini-0290:phyxminiproblems-problemphyxmini0290-timequantity}
  \lean{PhyXMiniProblems.ProblemPhyXMini0290.TimeQuantity}
  A signed time coordinate relative to the graph's time origin
\end{definition}
\begin{proof}
  This is the declared model definition of Time Quantity.
\end{proof}
\begin{definition}[Duration Quantity]
  \label{def:physics:phyx-mini-0290:phyxminiproblems-problemphyxmini0290-durationquantity}
  \lean{PhyXMiniProblems.ProblemPhyXMini0290.DurationQuantity}
  A positive-duration quantity, used here for the wave period
\end{definition}
\begin{proof}
  This is the declared model definition of Duration Quantity.
\end{proof}
\begin{definition}[Angular Frequency Quantity]
  \label{def:physics:phyx-mini-0290:phyxminiproblems-problemphyxmini0290-angularfrequencyquantity}
  \lean{PhyXMiniProblems.ProblemPhyXMini0290.AngularFrequencyQuantity}
  Angular frequency, with radians treated as dimensionless
\end{definition}
\begin{proof}
  This is the declared model definition of Angular Frequency Quantity.
\end{proof}
\begin{definition}[Wave Number Quantity]
  \label{def:physics:phyx-mini-0290:phyxminiproblems-problemphyxmini0290-wavenumberquantity}
  \lean{PhyXMiniProblems.ProblemPhyXMini0290.WaveNumberQuantity}
  Wave number, with radians treated as dimensionless
\end{definition}
\begin{proof}
  This is the declared model definition of Wave Number Quantity.
\end{proof}
\begin{definition}[Transverse Velocity Quantity]
  \label{def:physics:phyx-mini-0290:phyxminiproblems-problemphyxmini0290-transversevelocityquantity}
  \lean{PhyXMiniProblems.ProblemPhyXMini0290.TransverseVelocityQuantity}
  A signed transverse-velocity component
\end{definition}
\begin{proof}
  This is the declared model definition of Transverse Velocity Quantity.
\end{proof}
\begin{definition}[nonnegative Length Readout]
  \label{def:physics:phyx-mini-0290:phyxminiproblems-problemphyxmini0290-nonnegativelengthreadout}
  \lean{PhyXMiniProblems.ProblemPhyXMini0290.nonnegativeLengthReadout}
  Read a nonnegative length in the selected unit
\end{definition}
\begin{proof}
  This is the declared model definition of nonnegative Length Readout.
\end{proof}
\begin{definition}[signed Length Readout]
  \label{def:physics:phyx-mini-0290:phyxminiproblems-problemphyxmini0290-signedlengthreadout}
  \lean{PhyXMiniProblems.ProblemPhyXMini0290.signedLengthReadout}
  \uses{def:physics:phyx-mini-0290:phyxminiproblems-problemphyxmini0290-signedlengthquantity}
  Read a signed length coordinate in the selected unit
\end{definition}
\begin{proof}
  This is the declared model definition of signed Length Readout.
\end{proof}
\begin{definition}[time Readout]
  \label{def:physics:phyx-mini-0290:phyxminiproblems-problemphyxmini0290-timereadout}
  \lean{PhyXMiniProblems.ProblemPhyXMini0290.timeReadout}
  \uses{def:physics:phyx-mini-0290:phyxminiproblems-problemphyxmini0290-timequantity}
  Read a signed time coordinate in the selected unit
\end{definition}
\begin{proof}
  This is the declared model definition of time Readout.
\end{proof}
\begin{definition}[duration Readout]
  \label{def:physics:phyx-mini-0290:phyxminiproblems-problemphyxmini0290-durationreadout}
  \lean{PhyXMiniProblems.ProblemPhyXMini0290.durationReadout}
  \uses{def:physics:phyx-mini-0290:phyxminiproblems-problemphyxmini0290-durationquantity}
  Read a positive duration in the selected unit
\end{definition}
\begin{proof}
  This is the declared model definition of duration Readout.
\end{proof}
\begin{definition}[angular Frequency Readout]
  \label{def:physics:phyx-mini-0290:phyxminiproblems-problemphyxmini0290-angularfrequencyreadout}
  \lean{PhyXMiniProblems.ProblemPhyXMini0290.angularFrequencyReadout}
  \uses{def:physics:phyx-mini-0290:phyxminiproblems-problemphyxmini0290-angularfrequencyquantity}
  Read angular frequency in radians per selected time unit
\end{definition}
\begin{proof}
  This is the declared model definition of angular Frequency Readout.
\end{proof}
\begin{definition}[wave Number Readout]
  \label{def:physics:phyx-mini-0290:phyxminiproblems-problemphyxmini0290-wavenumberreadout}
  \lean{PhyXMiniProblems.ProblemPhyXMini0290.waveNumberReadout}
  \uses{def:physics:phyx-mini-0290:phyxminiproblems-problemphyxmini0290-wavenumberquantity}
  Read wave number in radians per selected length unit
\end{definition}
\begin{proof}
  This is the declared model definition of wave Number Readout.
\end{proof}
\begin{definition}[transverse Velocity Readout]
  \label{def:physics:phyx-mini-0290:phyxminiproblems-problemphyxmini0290-transversevelocityreadout}
  \lean{PhyXMiniProblems.ProblemPhyXMini0290.transverseVelocityReadout}
  \uses{def:physics:phyx-mini-0290:phyxminiproblems-problemphyxmini0290-transversevelocityquantity}
  Read transverse velocity in a coherent selected length/time unit system
\end{definition}
\begin{proof}
  This is the declared model definition of transverse Velocity Readout.
\end{proof}
\begin{definition}[String Wave Kind]
  \label{def:physics:phyx-mini-0290:phyxminiproblems-problemphyxmini0290-stringwavekind}
  \lean{PhyXMiniProblems.ProblemPhyXMini0290.StringWaveKind}
  The disturbance type described in the problem
\end{definition}
\begin{proof}
  This is the declared model definition of String Wave Kind.
\end{proof}
\begin{definition}[Propagation Direction]
  \label{def:physics:phyx-mini-0290:phyxminiproblems-problemphyxmini0290-propagationdirection}
  \lean{PhyXMiniProblems.ProblemPhyXMini0290.PropagationDirection}
  Propagation direction along the positive-x axis
\end{definition}
\begin{proof}
  This is the declared model definition of Propagation Direction.
\end{proof}
\begin{definition}[temporal Phase Sign]
  \label{def:physics:phyx-mini-0290:phyxminiproblems-problemphyxmini0290-temporalphasesign}
  \lean{PhyXMiniProblems.ProblemPhyXMini0290.temporalPhaseSign}
  \uses{def:physics:phyx-mini-0290:phyxminiproblems-problemphyxmini0290-propagationdirection}
  Sign multiplying the temporal phase in kx ± omega t + phi
\end{definition}
\begin{proof}
  This is the declared model definition of temporal Phase Sign.
\end{proof}
\begin{definition}[Graph Axis]
  \label{def:physics:phyx-mini-0290:phyxminiproblems-problemphyxmini0290-graphaxis}
  \lean{PhyXMiniProblems.ProblemPhyXMini0290.GraphAxis}
  The two labeled axes in the supplied graph
\end{definition}
\begin{proof}
  This is the declared model definition of Graph Axis.
\end{proof}
\begin{definition}[Axis Quantity]
  \label{def:physics:phyx-mini-0290:phyxminiproblems-problemphyxmini0290-axisquantity}
  \lean{PhyXMiniProblems.ProblemPhyXMini0290.AxisQuantity}
  Physical quantity printed beside each graph axis
\end{definition}
\begin{proof}
  This is the declared model definition of Axis Quantity.
\end{proof}
\begin{definition}[Wave Graph Landmark]
  \label{def:physics:phyx-mini-0290:phyxminiproblems-problemphyxmini0290-wavegraphlandmark}
  \lean{PhyXMiniProblems.ProblemPhyXMini0290.WaveGraphLandmark}
  Distinguished points of the plotted sinusoid, ordered from left to right
\end{definition}
\begin{proof}
  This is the declared model definition of Wave Graph Landmark.
\end{proof}
\begin{definition}[Displacement Time Graph]
  \label{def:physics:phyx-mini-0290:phyxminiproblems-problemphyxmini0290-displacementtimegraph}
  \lean{PhyXMiniProblems.ProblemPhyXMini0290.DisplacementTimeGraph}
  \uses{def:physics:phyx-mini-0290:phyxminiproblems-problemphyxmini0290-amplitudequantity, def:physics:phyx-mini-0290:phyxminiproblems-problemphyxmini0290-timequantity, def:physics:phyx-mini-0290:phyxminiproblems-problemphyxmini0290-graphaxis, def:physics:phyx-mini-0290:phyxminiproblems-problemphyxmini0290-axisquantity, def:physics:phyx-mini-0290:phyxminiproblems-problemphyxmini0290-wavegraphlandmark}
  Labels, units, and landmark coordinates supplied by the primary graph. landmarkTime is dimensionful; the numerical times are recorded separately in MatchesPrimaryWaveGraph rather than built into this structure
\end{definition}
\begin{proof}
  This is the declared model definition of Displacement Time Graph.
\end{proof}
\begin{definition}[Traveling String Wave Setup]
  \label{def:physics:phyx-mini-0290:phyxminiproblems-problemphyxmini0290-travelingstringwavesetup}
  \lean{PhyXMiniProblems.ProblemPhyXMini0290.TravelingStringWaveSetup}
  \uses{def:physics:phyx-mini-0290:phyxminiproblems-problemphyxmini0290-amplitudequantity, def:physics:phyx-mini-0290:phyxminiproblems-problemphyxmini0290-wavelengthquantity, def:physics:phyx-mini-0290:phyxminiproblems-problemphyxmini0290-signedlengthquantity, def:physics:phyx-mini-0290:phyxminiproblems-problemphyxmini0290-timequantity, def:physics:phyx-mini-0290:phyxminiproblems-problemphyxmini0290-durationquantity, def:physics:phyx-mini-0290:phyxminiproblems-problemphyxmini0290-angularfrequencyquantity, def:physics:phyx-mini-0290:phyxminiproblems-problemphyxmini0290-wavenumberquantity, def:physics:phyx-mini-0290:phyxminiproblems-problemphyxmini0290-transversevelocityquantity, def:physics:phyx-mini-0290:phyxminiproblems-problemphyxmini0290-stringwavekind, def:physics:phyx-mini-0290:phyxminiproblems-problemphyxmini0290-propagationdirection, def:physics:phyx-mini-0290:phyxminiproblems-problemphyxmini0290-displacementtimegraph}
  Independent physical quantities and observables for the string wave. Displacement and transverse velocity remain independent dimensionful fields; the governing-law premise below relates them to the wave parameters
\end{definition}
\begin{proof}
  This is the declared model definition of Traveling String Wave Setup.
\end{proof}
\begin{definition}[Matches String Wave Problem Data]
  \label{def:physics:phyx-mini-0290:phyxminiproblems-problemphyxmini0290-matchesstringwaveproblemdata}
  \lean{PhyXMiniProblems.ProblemPhyXMini0290.MatchesStringWaveProblemData}
  \uses{def:physics:phyx-mini-0290:phyxminiproblems-problemphyxmini0290-nonnegativelengthreadout, def:physics:phyx-mini-0290:phyxminiproblems-problemphyxmini0290-signedlengthreadout, def:physics:phyx-mini-0290:phyxminiproblems-problemphyxmini0290-timereadout, def:physics:phyx-mini-0290:phyxminiproblems-problemphyxmini0290-travelingstringwavesetup}
  Problem-statement data: the disturbance is a right-moving transverse sinusoid of wavelength 20 cm, and the requested event is x = 0, t = 5.0 s. No transverse-velocity value occurs in these premises
\end{definition}
\begin{proof}
  This is the declared model definition of Matches String Wave Problem Data.
\end{proof}
\begin{definition}[Matches Primary Wave Graph]
  \label{def:physics:phyx-mini-0290:phyxminiproblems-problemphyxmini0290-matchesprimarywavegraph}
  \lean{PhyXMiniProblems.ProblemPhyXMini0290.MatchesPrimaryWaveGraph}
  \uses{def:physics:phyx-mini-0290:phyxminiproblems-problemphyxmini0290-nonnegativelengthreadout, def:physics:phyx-mini-0290:phyxminiproblems-problemphyxmini0290-signedlengthreadout, def:physics:phyx-mini-0290:phyxminiproblems-problemphyxmini0290-timereadout, def:physics:phyx-mini-0290:phyxminiproblems-problemphyxmini0290-durationreadout, def:physics:phyx-mini-0290:phyxminiproblems-problemphyxmini0290-transversevelocityreadout, def:physics:phyx-mini-0290:phyxminiproblems-problemphyxmini0290-travelingstringwavesetup}
  Calibrated evidence from the supplied bitmap. The five equal horizontal grid intervals run from 0 s to the second crest at 12.5 s; the printed 10 is at the ascending zero. Extrema touch ±y\_s, with y\_s = 4 cm. The slope signs record what is visibly rising or descending, but give no numerical velocity answer
\end{definition}
\begin{proof}
  This is the declared model definition of Matches Primary Wave Graph.
\end{proof}
\begin{definition}[Has Physical Traveling Wave Parameters]
  \label{def:physics:phyx-mini-0290:phyxminiproblems-problemphyxmini0290-hasphysicaltravelingwaveparameters}
  \lean{PhyXMiniProblems.ProblemPhyXMini0290.HasPhysicalTravelingWaveParameters}
  \uses{def:physics:phyx-mini-0290:phyxminiproblems-problemphyxmini0290-nonnegativelengthreadout, def:physics:phyx-mini-0290:phyxminiproblems-problemphyxmini0290-durationreadout, def:physics:phyx-mini-0290:phyxminiproblems-problemphyxmini0290-angularfrequencyreadout, def:physics:phyx-mini-0290:phyxminiproblems-problemphyxmini0290-wavenumberreadout, def:physics:phyx-mini-0290:phyxminiproblems-problemphyxmini0290-travelingstringwavesetup}
  Positivity and the conventional principal representative for phase. The phase range does not fix its value; the graph and governing laws select it
\end{definition}
\begin{proof}
  This is the declared model definition of Has Physical Traveling Wave Parameters.
\end{proof}
\begin{definition}[Satisfies Sinusoidal Traveling Wave Laws]
  \label{def:physics:phyx-mini-0290:phyxminiproblems-problemphyxmini0290-satisfiessinusoidaltravelingwavelaws}
  \lean{PhyXMiniProblems.ProblemPhyXMini0290.SatisfiesSinusoidalTravelingWaveLaws}
  \uses{def:physics:phyx-mini-0290:phyxminiproblems-problemphyxmini0290-signedlengthquantity, def:physics:phyx-mini-0290:phyxminiproblems-problemphyxmini0290-timequantity, def:physics:phyx-mini-0290:phyxminiproblems-problemphyxmini0290-nonnegativelengthreadout, def:physics:phyx-mini-0290:phyxminiproblems-problemphyxmini0290-signedlengthreadout, def:physics:phyx-mini-0290:phyxminiproblems-problemphyxmini0290-timereadout, def:physics:phyx-mini-0290:phyxminiproblems-problemphyxmini0290-durationreadout, def:physics:phyx-mini-0290:phyxminiproblems-problemphyxmini0290-angularfrequencyreadout, def:physics:phyx-mini-0290:phyxminiproblems-problemphyxmini0290-wavenumberreadout, def:physics:phyx-mini-0290:phyxminiproblems-problemphyxmini0290-transversevelocityreadout, def:physics:phyx-mini-0290:phyxminiproblems-problemphyxmini0290-temporalphasesign, def:physics:phyx-mini-0290:phyxminiproblems-problemphyxmini0290-travelingstringwavesetup}
  Generic laws for a sinusoidal traveling string wave: k lambda = 2 pi; omega T = 2 pi; y(x,t) = y\_m sin(kx ± omega t + phi), with the minus sign for travel along positive x; transverse particle velocity is the time derivative of that displacement. All equations use coherent named-unit readouts. The velocity law is the dimensionful observable interface corresponding to Mathlib's scalar HasDerivAt.sin chain rule. None of these laws mentions the requested event or a displayed answer value
\end{definition}
\begin{proof}
  This is the declared model definition of Satisfies Sinusoidal Traveling Wave Laws.
\end{proof}
\begin{definition}[Answer Choice]
  \label{def:physics:phyx-mini-0290:phyxminiproblems-problemphyxmini0290-answerchoice}
  \lean{PhyXMiniProblems.ProblemPhyXMini0290.AnswerChoice}
  Labels of the four velocity choices printed with the problem
\end{definition}
\begin{proof}
  This is the declared model definition of Answer Choice.
\end{proof}
\begin{definition}[displayed Velocity In Centimeters Per Second]
  \label{def:physics:phyx-mini-0290:phyxminiproblems-problemphyxmini0290-displayedvelocityincentimeterspersecond}
  \lean{PhyXMiniProblems.ProblemPhyXMini0290.displayedVelocityInCentimetersPerSecond}
  \uses{def:physics:phyx-mini-0290:phyxminiproblems-problemphyxmini0290-answerchoice}
  Displayed transverse velocity in centimeters per second
\end{definition}
\begin{proof}
  This is the declared model definition of displayed Velocity In Centimeters Per Second.
\end{proof}
\begin{definition}[recorded Dataset Answer]
  \label{def:physics:phyx-mini-0290:phyxminiproblems-problemphyxmini0290-recordeddatasetanswer}
  \lean{PhyXMiniProblems.ProblemPhyXMini0290.recordedDatasetAnswer}
  \uses{def:physics:phyx-mini-0290:phyxminiproblems-problemphyxmini0290-answerchoice}
  The answer label recorded in the supplied dataset metadata
\end{definition}
\begin{proof}
  This is the declared model definition of recorded Dataset Answer.
\end{proof}
\begin{definition}[Matches Displayed Velocity]
  \label{def:physics:phyx-mini-0290:phyxminiproblems-problemphyxmini0290-matchesdisplayedvelocity}
  \lean{PhyXMiniProblems.ProblemPhyXMini0290.MatchesDisplayedVelocity}
  \uses{def:physics:phyx-mini-0290:phyxminiproblems-problemphyxmini0290-answerchoice, def:physics:phyx-mini-0290:phyxminiproblems-problemphyxmini0290-displayedvelocityincentimeterspersecond}
  Agreement with a one-decimal-place displayed velocity
\end{definition}
\begin{proof}
  This is the declared model definition of Matches Displayed Velocity.
\end{proof}
\begin{lemma}[angular Frequency In Radians Per Second eq pi over five]
  \label{lem:physics:phyx-mini-0290:phyxminiproblems-problemphyxmini0290-angularfrequencyinradianspersecond-eq-pi-over-five}
  \lean{PhyXMiniProblems.ProblemPhyXMini0290.angularFrequencyInRadiansPerSecond_eq_pi_over_five}
  \uses{def:physics:phyx-mini-0290:phyxminiproblems-problemphyxmini0290-angularfrequencyreadout, def:physics:phyx-mini-0290:phyxminiproblems-problemphyxmini0290-travelingstringwavesetup, def:physics:phyx-mini-0290:phyxminiproblems-problemphyxmini0290-matchesprimarywavegraph, def:physics:phyx-mini-0290:phyxminiproblems-problemphyxmini0290-satisfiessinusoidaltravelingwavelaws}
  The 10 s period gives omega = pi/5 rad/s
\end{lemma}
\begin{proof}
  The conclusion follows from the governing relations and figure readouts named in the statement of angular Frequency In Radians Per Second eq pi over five.
\end{proof}
\begin{lemma}[wave Number In Radians Per Centimeter eq pi over ten]
  \label{lem:physics:phyx-mini-0290:phyxminiproblems-problemphyxmini0290-wavenumberinradianspercentimeter-eq-pi-over-ten}
  \lean{PhyXMiniProblems.ProblemPhyXMini0290.waveNumberInRadiansPerCentimeter_eq_pi_over_ten}
  \uses{def:physics:phyx-mini-0290:phyxminiproblems-problemphyxmini0290-wavenumberreadout, def:physics:phyx-mini-0290:phyxminiproblems-problemphyxmini0290-travelingstringwavesetup, def:physics:phyx-mini-0290:phyxminiproblems-problemphyxmini0290-matchesstringwaveproblemdata, def:physics:phyx-mini-0290:phyxminiproblems-problemphyxmini0290-satisfiessinusoidaltravelingwavelaws}
  The stated 20 cm wavelength gives k = pi/10 rad/cm. This records the role of wavelength even though the requested particle velocity depends only on amplitude, angular frequency, and phase at the event
\end{lemma}
\begin{proof}
  The conclusion follows from the governing relations and figure readouts named in the statement of wave Number In Radians Per Centimeter eq pi over ten.
\end{proof}
\begin{lemma}[phase Offset Radians eq pi]
  \label{lem:physics:phyx-mini-0290:phyxminiproblems-problemphyxmini0290-phaseoffsetradians-eq-pi}
  \lean{PhyXMiniProblems.ProblemPhyXMini0290.phaseOffsetRadians_eq_pi}
  \uses{def:physics:phyx-mini-0290:phyxminiproblems-problemphyxmini0290-travelingstringwavesetup, def:physics:phyx-mini-0290:phyxminiproblems-problemphyxmini0290-matchesstringwaveproblemdata, def:physics:phyx-mini-0290:phyxminiproblems-problemphyxmini0290-matchesprimarywavegraph, def:physics:phyx-mini-0290:phyxminiproblems-problemphyxmini0290-hasphysicaltravelingwaveparameters, def:physics:phyx-mini-0290:phyxminiproblems-problemphyxmini0290-satisfiessinusoidaltravelingwavelaws}
  At the graph origin the displacement vanishes and the curve rises. For a right-moving wave written with the problem's sine convention and principal phase representative, this selects phi = pi
\end{lemma}
\begin{proof}
  The conclusion follows from the governing relations and figure readouts named in the statement of phase Offset Radians eq pi.
\end{proof}
\begin{lemma}[transverse Velocity At Five Seconds exact]
  \label{lem:physics:phyx-mini-0290:phyxminiproblems-problemphyxmini0290-transversevelocityatfiveseconds-exact}
  \lean{PhyXMiniProblems.ProblemPhyXMini0290.transverseVelocityAtFiveSeconds_exact}
  \uses{def:physics:phyx-mini-0290:phyxminiproblems-problemphyxmini0290-transversevelocityreadout, def:physics:phyx-mini-0290:phyxminiproblems-problemphyxmini0290-travelingstringwavesetup, def:physics:phyx-mini-0290:phyxminiproblems-problemphyxmini0290-matchesstringwaveproblemdata, def:physics:phyx-mini-0290:phyxminiproblems-problemphyxmini0290-matchesprimarywavegraph, def:physics:phyx-mini-0290:phyxminiproblems-problemphyxmini0290-hasphysicaltravelingwaveparameters, def:physics:phyx-mini-0290:phyxminiproblems-problemphyxmini0290-satisfiessinusoidaltravelingwavelaws}
  At x = 0, t = 5 s, the phase is zero modulo a full turn and the exact model velocity is -4 pi / 5 cm/s
\end{lemma}
\begin{proof}
  The conclusion follows from the governing relations and figure readouts named in the statement of transverse Velocity At Five Seconds exact.
\end{proof}
% --- Archon declaration topology end ---
% --- Archon physics formalization source end ---
